\section{Thiết kế khối nguồn}
Khối nguồn là nền tảng cho toàn bộ node. Theo sơ đồ nguyên lý, mạch có đầu vào nguồn ngoài qua jack DC và header cấp nguồn. Từ nguồn đầu vào, mạch tạo ra các rail điện áp phục vụ từng nhóm linh kiện: điện áp cao hơn cho tải hoặc relay, 5V cho một số cảm biến và mạch điều khiển, 3.3V cho module LoRa và các linh kiện logic mức thấp.

\subsection{Nguồn đầu vào và bảo vệ}
Đầu vào nguồn được thiết kế với jack DC và diode bảo vệ. Diode giúp hạn chế rủi ro khi người dùng đấu nhầm cực hoặc có xung ngược từ tải. Trong môi trường thử nghiệm, việc cắm nhầm adapter hoặc đảo cực dây cấp nguồn là lỗi phổ biến, vì vậy phần bảo vệ nguồn là cần thiết. Ngoài diode bảo vệ, mạch còn bố trí tụ lọc ở các rail nguồn để giảm dao động điện áp khi relay đóng cắt hoặc khi module LoRa phát.

Một node IoT có nhiều loại tải khác nhau. Cảm biến thường tiêu thụ dòng nhỏ và yêu cầu nguồn ổn định. Relay và LED hồng ngoại có dòng tức thời cao hơn. LoRa khi phát cũng tạo dòng đỉnh. Nếu nguồn không đủ ổn định, node có thể reset, giá trị ADC dao động hoặc mất gói LoRa. Do đó khối nguồn cần được kiểm tra trước khi đánh giá các chức năng khác.

\begin{figure}[h!]
    \centering
    \begin{tikzpicture}[node distance=1.5cm, every node/.style={font=\small}]
        \node[draw, minimum width=2.5cm, minimum height=0.9cm] (vin) {DC Jack / 24V};
        \node[draw, right=of vin, minimum width=2.4cm, minimum height=0.9cm] (protect) {Diode bảo vệ};
        \node[draw, right=of protect, minimum width=2.4cm, minimum height=0.9cm] (buck) {Buck 5V};
        \node[draw, right=of buck, minimum width=2.4cm, minimum height=0.9cm] (ldo) {AMS1117 3.3V};
        \node[draw, below=of buck, minimum width=2.4cm, minimum height=0.9cm] (relay) {Relay / cảm biến};
        \node[draw, below=of ldo, minimum width=2.4cm, minimum height=0.9cm] (lora) {LoRa RA-02};
        \draw[->, thick] (vin) -- (protect);
        \draw[->, thick] (protect) -- (buck);
        \draw[->, thick] (buck) -- (ldo);
        \draw[->, thick] (buck) -- (relay);
        \draw[->, thick] (ldo) -- (lora);
    \end{tikzpicture}
    \caption{Luồng cấp nguồn chính của node}
    \label{fig:power_flow}
\end{figure}

\subsection{Mạch hạ áp 5V}
Rail 5V cấp cho các cảm biến phổ biến, relay, transistor điều khiển và một số module ngoại vi. Nếu sử dụng nguồn đầu vào 24V, việc hạ áp bằng tuyến tính sẽ gây tổn hao nhiệt lớn. Vì vậy sơ đồ sử dụng khối hạ áp dạng switching để tạo 5V hiệu quả hơn. Tụ đầu vào và đầu ra của khối hạ áp có nhiệm vụ giảm ripple và ổn định điện áp.

Đối với hệ thống có relay, dòng tải thay đổi theo trạng thái đóng cắt. Mỗi khi relay hút, cuộn dây tạo dòng đột biến; khi relay nhả, năng lượng trong cuộn dây có thể gây xung ngược. Mạch cần có diode dập xung tại cuộn relay và tụ lọc đủ gần rail 5V để tránh ảnh hưởng đến vi điều khiển.

\subsection{Mạch ổn áp 3.3V}
Module LoRa RA-02 hoạt động ở mức 3.3V. Sơ đồ sử dụng AMS1117-3.3V để tạo rail 3.3V từ 5V. AMS1117 là ổn áp tuyến tính dễ dùng, phù hợp với dòng tiêu thụ mức vừa phải. Tuy nhiên cần lưu ý nhiệt lượng nếu dòng 3.3V tăng cao. Trong thiết kế hiện tại, rail 3.3V chủ yếu cấp cho LoRa nên phương án này phù hợp.

Việc tách 5V và 3.3V giúp bảo vệ module LoRa khỏi mức logic không phù hợp. Khi Arduino Pro Mini chạy 5V, các chân SPI có thể tạo mức logic 5V. Vì RA-02 chỉ chịu 3.3V, thiết kế thực tế cần bảo đảm mức logic phù hợp, có thể thông qua lựa chọn Arduino Pro Mini 3.3V hoặc bổ sung mạch chuyển mức logic. Đây là điểm cần kiểm tra kỹ khi lắp mạch thật.

\subsection{LED báo nguồn}
Sơ đồ có các LED báo nguồn cho các rail như 5V, nguồn LoRa, nguồn module và LED trạng thái. LED báo nguồn giúp quá trình kiểm tra nhanh hơn: chỉ cần quan sát có thể biết rail nào đang hoạt động. Tuy nhiên điện trở hạn dòng cần được chọn phù hợp để tránh tiêu thụ dòng không cần thiết, đặc biệt nếu hệ thống hướng đến vận hành lâu dài bằng nguồn pin.

\begin{longtable}{|p{3cm}|p{4cm}|p{5.8cm}|}
\hline
\textbf{Rail nguồn} & \textbf{Phụ tải chính} & \textbf{Lưu ý kiểm thử} \\
\hline
24V hoặc nguồn vào & Adapter, đầu cấp ngoài, tải mở rộng & Kiểm tra cực tính, điện áp không tải và điện áp khi relay đóng. \\
\hline
5V & Cảm biến, relay, transistor điều khiển, Arduino nếu bản 5V & Đo ripple khi relay hút và khi LoRa phát. \\
\hline
3.3V & LoRa RA-02, logic mức thấp & Không cấp nhầm 5V cho module LoRa; kiểm tra dòng đỉnh khi truyền. \\
\hline
GND & Toàn bộ hệ thống & Cần nối mass chung giữa nguồn, Arduino, cảm biến, relay và LoRa. \\
\hline
\end{longtable}

\subsection{Rủi ro và biện pháp giảm nhiễu}
Các rủi ro chính của khối nguồn gồm sụt áp khi relay đóng, nhiễu ảnh hưởng đến ADC, nóng ổn áp 3.3V và đấu nhầm nguồn. Biện pháp giảm rủi ro gồm dùng adapter đủ dòng, bố trí tụ lọc gần các module tiêu thụ dòng xung, đi dây mass hợp lý, tách đường công suất và đường tín hiệu analog, đồng thời kiểm tra điện áp từng rail trước khi gắn module nhạy như LoRa hoặc cảm biến pH.

Trong quá trình thử nghiệm, nên đo nguồn ở ba trạng thái: node chỉ đọc cảm biến, node truyền LoRa và node vừa truyền LoRa vừa đóng relay. Nếu điện áp không ổn định ở trạng thái thứ ba, cần tăng khả năng cấp dòng hoặc bổ sung tụ lọc trước khi đánh giá phần mềm.
